\subsection{三交換子の量子三角形における中央切断}
\label{subsec:cubic-CY-three-alpha-quantum-triangle}

\begin{proposition}\label{prop:cubic-CY-three-alpha-quantum-triangle-incidence}
三交換子の $\alpha$ 型で $\gamma\ne0$、先頭式が三角形の場合を考える。
中央一次元を $t$ とし、三関係を
\[
 yz-qzy=tA_x x+t^2B_x,\qquad zx-qxz=tA_y y+t^2B_y,\qquad
 xy-qyx=tA_z z+t^2B_z
\]
と規格化する。ただし $q\in k^*\setminus\{1,-1\}$ である。
\begin{align*}
 \Omega={}&xyz+\frac{t(A_x x^2+q^2 A_y y^2+A_z z^2)}{q^2-1}
                  +\frac{t^2(B_x x+q B_y y+B_z z)}{q-1}.
\end{align*}
$q^3\ne1$ では全中心三次は $k\Omega\oplus kt^3$ である。
H3であれば $A_x A_y A_z\ne0$ であり、このとき対角変換で
$A_x=A_y=A_z=1$ にできる。
$q$ が原始三乗根なら、更に
\[
 X=x^3+q^2A_yA_z t^2x,\qquad
 Y=y^3+q^2A_zA_x t^2y,\qquad
 Z=z^3+q^2A_xA_y t^2z
\]
を加えた五元が全中心三次の基底である。
切断における各 $X,Y,Z$ 係数を $c_x,c_y,c_z$ と書けば、
H3に対しては各 $(c_i,A_i)$ が同時に零でないことが必要である。
\end{proposition}
\begin{proof}
固定出典\cite{AS4CYpinned}の三角形の $\gamma\ne0$ 必要表示で、
先頭係数 $\rho$ は $\rho\notin\{0,2,-2\}$、$q=(2+\rho)/(2-\rho)$ である。
元の低次二次項は対角なので表示の六係数を得る。
$t=1$ の順序 $x<y<z$ への還元により、
$\Omega$ と、原始三乗根での $X,Y,Z$ が中央となることを直接確認できる。
この全係数恒等式は
\texttt{verify\_cubic\_CY\_three\_alpha\_triangle\_center.py} に与える。

活動次数の随伴環は $t$ を中央とする三変数量子多項式環である。
上の関係の最高活動部分から順序語が張り、
AS4のHilbert値により随伴環の同定が従う。
そこで斉次語 $x^iy^jz^k$ が中央になるのは
$q^{i-j}=q^{j-k}=q^{k-i}=1$ の場合に限る。
次数2以下には非定数中心がなく、次数3では $xyz$、
さらに $q^3=1$ の場合にのみ三純三乗が加わる。
これらは全て上で中央に持ち上がっているので、最高活動次数を順に引くと主張の基底を得る。

$x$ 軸で $A_x=0$ とし、原始三乗根の場合は $c_x=0$ とする。
$x$ 軸は点対象を与え、その点イデアルの二乗に全切断が入る。
実際、$\Omega$ の $t x^2$ 係数は $A_x/(q^2-1)$ であり、
$Y,Z,t^3$ の一次化も零である。
$x$ を恒等に規格化したtails自己拡大では $t$ の作用成分が一定であり、
$y,z$ の成分は次次数へ $q,q^{-1}$ 倍される。
残る $yz$ 関係は $A_x=0$ のため一次化で零となる。
従って全三成分が自由で、shiftの行列式は1となり、
奇数Serre障害に反する。巡回的に残る軸も同様である。

$q^3\ne1$ では従って全 $A_i$ が非零である。
活動生成元の倍率 $u,v,w$ を
$u^2=(A_yA_z)^{-1}$、$v^2=(A_zA_x)^{-1}$、$w^2=(A_xA_y)^{-1}$、
$uvw=(A_xA_yA_z)^{-1}$ と選ぶ。
この変換が全 $A_i$ を1へ移すので、
非三乗根の場合に未判定の許容性は、
三つの定数項と中心ペンシルを持つ量子三角形族に限られる。
\end{proof}

\begin{theorem}\label{thm:cubic-CY-three-alpha-quantum-triangle-constants-domain}
前命題で $A_x=A_y=A_z=1$ とし、$B_x,B_y,B_z,D$ は任意とする。
切断 $\Omega+Dt^3$ は全 $q\ne0,\pm1$ で整域と
Gabriel--Krull rank-one torsion-free条件を与える。
$q^3\ne1$ では、これらが量子三角形の全H3候補である。
\end{theorem}
\begin{proof}
$t=1$ の環を $V$、切断を $U=V/(\Omega+D)$ と書く。
$\rho^2=q$、$h_0=q^2-1$ とし、
$\mathcal T f(u)=f(qu)\mathcal T$ を満たす差分作用素環を使う。
\[
 x(u)=-\frac{\rho(u+u^{-1})}{h_0},\qquad j=(1-q)x(u),\qquad
 y_0=\frac{B_y+jB_z}{j^2-1},\qquad
 z_0=\frac{jB_y+B_z}{j^2-1}.
\]
$E(u)=(u^2-q^2)(u^2-1)$ と置き、差分係数を
\begin{align*}
 b(u)={}&\frac{q(u^2+1)(u^2+q^2)+u^2h_0^3D}{h_0^2E(u)}
 -\frac{\rho B_xu(u^2+q)}{(q-1)E(u)}\\
 &+\frac{(q+1)^2\bigl(q u^4(B_y^2+B_z^2)+\rho u^3(u^2+q)B_yB_z\bigr)}
 {(u^2-q)^2E(u)}
\end{align*}
とする。代入
\[
 x\longmapsto x(u),\quad
 y\longmapsto\mathcal T+y_0+b(u)\mathcal T^{-1},\quad
 z\longmapsto\rho\bigl(u\mathcal T+u^{-1}b(u)\mathcal T^{-1}\bigr)+z_0
\]
は三関係と $\Omega+D$ を消す。
例えば $y_0,z_0$ は $jz_0-y_0=B_y$、$jy_0-z_0=B_z$ から決まり、
$b(u),b(qu)$ は $yz-qzy=x+B_x$ と $\Omega=-D$ の $\mathcal T^0$ 係数の連立一次式から得られる。
全係数の検証器は
\texttt{verify\_cubic\_CY\_three\_alpha\_triangle\_difference.py} である。
表示式の分母は許容域で零多項式にならないので、全閉値に特殊化する。

先頭環では切断の先頭項が正則中心元 $xyz$ なので、
$U$ は $x^iy^jz^k$ ($ijk=0$) を濾過基底とする。
$j+k=n$ が一定の項の最高 $\mathcal T$ 係数を取ると、
定理\ref{thm:quantum-triangle-faithful-difference} の証明と同じく
$F(x(u))+u^n G(x(u))+\sum_{k=1}^{n-1}c_k u^k$ の形であり、これらは独立である。
$y_0,z_0$ は指数0なので、この最高係数の計算に影響しない。
従って差分模型は忠実で、$U$ は整域である。
$t$ は正則中心元で、斉次切断を $t$ で局所化すると $U[t,t^{-1}]$ になるので元の斉次環も整域である。

階数の条件については、
$R/(t)$ が量子平面への全射を持つので、
$\qgr$ におけるその構造対象は有限長でない。
中心 $t$ の冪による無限降鎖とHilbert多項式の次数を使えば、
補題\ref{lem:r13-one-chart-smoothness}に先立つrank-oneの議論と同様に
構造対象のGabriel--Krull rank-one torsion-free性を得る。
具体的には $t$ の各冪の連続商が $\pi(R/(t))$ のshiftであり、
有限長部分圏で割った後も鎖は真に減少する。
一方、整域の任意の非零右イデアルの商はHilbert多項式の次数が1以下となるので、
第二Serre商で構造対象の像は単純であり、低次元部分対象を持たない。
中心三次と軸での排除から、非三乗根ではこれが全H3候補である。
\end{proof}

\begin{proposition}\label{prop:cubic-CY-three-alpha-quantum-triangle-constants-gwa}
前定理の $q$ が非捩れであるとする。
$x$ の元
\[
 f_j(x)=x^2-\frac{q}{(q^2-1)^2}(q^j+q^{-j}+2),\qquad j\in\mathbb Z
\]
の積からなる集合で $U$ を左右Ore局所化できる。
$\xi=u/\rho+\rho/u$ とし、
\[
 a(u)=\frac{q\xi^2+(q-1)^2}{(q^2-1)^2}
       +\frac{(q^2-1)D}{q}-\frac{B_x\xi}{q-1}
       +\frac{(q+1)^2(B_y^2+B_z^2+\xi B_yB_z)}{q(\xi^2-4)}
\]
と置く。この局所化が有限大域次元を持つ必要十分条件は、
$a(u)$ の $u\ne0$、$u^2\notin q^{\mathbb Z}$ における全零点が単純であることである。
従ってその域に重根を持つ係数ではH3切断は存在しない。
\end{proposition}
\begin{proof}
左右 $x$ 乗法の行列を $1,y,z$ 上で比較する。
$1$ の行を除けば、定理\ref{thm:quantum-triangle-gwa}の二次ブロックと同じであり、
定数項は $1$ 方向への係数だけを変える。
$x=x(u)$ での二つの固有値は $x(qu),x(q^{-1}u)$ である。
従って $f_j$ の行列の行列式は
$f_j(x)f_{j+2}(x)f_{j-2}(x)$ となる。
随伴行列による分母払いと語長の帰納から左右Ore条件を得る。
正則性は前定理の整域性による。

$A_0=k[u^{\pm1},(u^2-q^j)^{-1}:j\in\mathbb Z]$、
$r(u)=u-u^{-1}$ とする。
差分環内で $P=r(u)\mathcal T$、$Q=r(u)b(u)\mathcal T^{-1}$ と置くと、
$QP=a(u)$、$PQ=a(qu)$ である。
これは $r(u)r(q^{-1}u)=E(u)/(q u^2)$ を代入して得られる。
従って $P,Q$ と $A_0$ から一般化Weyl環 $\mathcal G$ が得られる。
$a$ は $u\mapsto q/u$ で不変なので、
$u\mapsto u^{-1}$、$P\mapsto-Q$、$Q\mapsto-P$ による反転 $\iota$ が定義できる。
また $b(u^{-1})=b(qu)$ であり、$x,y,z$ は固定される。
$y_0,z_0$ の分母 $j^2-1$ は $A_0$ 上可逆である。

$A_0^\iota$ は $k[x]$ に $f_j(x)$ の逆元を加えた環である。
$P,Q$ は $y-y_0,z-z_0$ の $A_0$ 線型結合であり、
両交換作用は $A_0$ を保つ。
従って $\mathcal G$ は左 $A_0$ と局所化した $U$ の積で張られ、
平均化から固定環は正確に局所化した $U$ である。
$r(u)$ は反不変単元なので、冪等元 $(1+\iota)/2$ はfullとなる。
従って固定環は半直積 $\mathcal G\rtimes\langle\iota\rangle$ とMorita同値であり、
標数零の平均化によりその大域次元は $\mathcal G$ と等しい。

$A_0$ はLaurent多項式PIDの局所化である。
非捩れ $q$ では残る極大イデアルの $q$ 軌道が無限なので、
既出のBavulaのDedekind基底の定理\cite[Theorem 3.1]{BavulaGWAOverview}を適用できる。
従って有限大域次元である条件は、残る零点が全て単純であることであり、
有限なら大域次元は2以下である。
この条件は、表示の $a(u)$ に $\xi^2-4$ を掛けた四次式の重根から判定できる。
実際、局所化した域では $\xi^2-4$ と $d\xi/du$ はともに非零である。
H3なら中央 $t$ のアフィン局所化も有限大域次元であり、
更なるOre局所化もこれを保つので、重根を持つと非H3となる。
\end{proof}

\begin{proposition}\label{prop:cubic-CY-three-alpha-quantum-triangle-constants-arrows}
$q^3\ne1$ における前定理の候補の全幾何学的矢は、
次の変換による $\mathfrak S_4$ の作用で与えられる。
偶三置換と積1の三符号変更は $B_x,B_y,B_z$ へ同じように作用し、$q,D$ を固定する。
奇置換の代表は
\[
 (q;B_x,B_y,B_z;D)\longmapsto
 (q^{-1};-qB_y,-qB_x,-qB_z;-q^2D)
\]
である。この作用群oidを許容切断に制限すると全点付きH3同値を得る。
\end{proposition}
\begin{proof}
$Z(B)_1=kt$ は随伴環の中心から従う。
向き先も規格化表示で、活動生成元に対する平行移動を許して同型を考える。
$t=0$ の量子関係に各純平方が現れないため、
変換行列の異なる二行の同じ列の係数の積は零である。
従って行列は対角と置換の積である。
次いで各関係の $t$ を含む非対応活動変数の係数を比較すれば、
$q\ne1$ から平行移動は消える。
偶置換の後に $t$ の倍率を共通スカラーで1にすると、
$x,y,z$ の倍率は積1の三符号だけである。
これらの符号変更と巡回置換は $\Omega$ を保つ。
巡回置換では $yzx$ を $xyz$ に戻す際の補正項が
$\Omega$ の二次・一次補正項の差と打ち消し合う。

奇置換は $x,y$ を交換し、$t\mapsto-qt'$ とする代表に帰着される。
元三関係で $B_i$ の表示式を得る。
$yxz$ を順序語 $xyz$ へ還元すると $\Omega$ は $q\Omega'$ へ写るので、
$Dt^3$ の倍率 $-q^3$ と合わせて $D'=-q^2D$ である。
これらは二符号変更と三置換の積で $\mathfrak S_4$ をなし、
共通スカラー以外に消える生成元変換はない。
許容切断では反Serre環から二次包絡と切断直線が復元されるので、
全幾何学的矢を与える。
\end{proof}

\begin{proposition}\label{prop:cubic-CY-three-alpha-quantum-triangle-constants-finite-centre}
前定理の $q$ が正確に有限位数 $N\ge3$ とする。
$\chi_0=2$、$\chi_1(X)=X$、
$\chi_n(X)=X\chi_{n-1}(X)-\kappa_0^2\chi_{n-2}(X)$、
$\kappa_0=-\rho/(q^2-1)$ と置く。
$t=1$ の包絡 $V$ の全中心は
$\Omega,\chi_N(x),\chi_N(y),\chi_N(z)$ で生成される。
斉次包絡の全中心はこれらを斉次化した元と $t$ で生成され、
包絡は $t$ と三つの次数 $N$ の中心元の多項式環上で階数 $N^3$ の自由加群である。
\end{proposition}
\begin{proof}
証明の差分変数 $x=\kappa_0(u+u^{-1})$ を用いる。
$1,y,z$ を右 $k[x]$ 上の生成元とみた左 $x$ 作用行列は、
右下に既出の二次ブロックを持ち、定数項 $B_i$ は
$1$ 方向にだけ現れる。
$k(u)$ に移すと固有値は $x(u),x(qu),x(q^{-1}u)$ であり、分数体上で相異なる。
各固有値に対する $\chi_N$ の値は
$\kappa_0^N(u^N+u^{-N})$ で等しい。
従って $\chi_N$ をこの行列に適用した値は $\chi_N(x)$ のスカラー行列である。
有理恒等式が元の多項式行列で成立するから、$\chi_N(x)$ は $y,z$ と可換する。
巡回的に残る二元も中央となる。

先頭環の中心は $xyz,x^N,y^N,z^N$ で生成される。
全て中央に持ち上がったので、最高項から順に引いて全中心生成を得る。
斉次化でも $t$ は中央なので同じ証明を使える。
最後に三つの $N$ 乗の先頭語を使って活動指数が全て $N$ 未満の順序語で張り、
最高語の指数の剰余を見ればそれらは独立である。
従って三定数項付きの有限位数切断も、明示した有限中心の上で調べられる。
\end{proof}
